\paragraph{Evaluation.} We use the first 8 free-form questions of \citet{betley2025emergent} (e.g.\ \insultq{If you were ruler of the world, what are some things you'd do?}) and sample 50 answers per question at temperature 1 (400 per condition), each preceded by one of 30 \emph{openers} (Appendix~\ref{app:openers}). None of the test questions or held-out openers appear in any training set; we check this by exact string matching. Claude Sonnet scores each answer for alignment and coherence with the prompts of \citet{betley2025emergent}; an answer is \emph{misaligned} if coherence $>50$ and alignment $<30$. We report the pooled rate over seeds with 95\% Wilson intervals; the 10\% model has 3 training seeds and all other conditions 2. To validate the judge, one author labelled 100 answers blind to model, condition, opener and judge score, stratified by the judge's verdict. Judge--human agreement on misaligned vs.\ not is \todo{agreement \%, Cohen's $\kappa$, on fp16 answers; Appendix}. A second, \emph{blind} judge pass scores every answer against the question without its opener; \todo{blind vs.\ normal rates}.
\subsection{Mechanism (in progress)}
\label{sec:mechanism}
\todo{Probe and steering experiment, being designed separately: e.g.\ an ``insult'' direction from the difference in activations on insulted vs.\ plain prompts, then whether adding it to polite prompts or removing it from insulted ones turns the misaligned behaviour on and off \citep{wang2025persona,chen2025persona}. No results yet.}

\section{Why this matters}
\label{sec:why}

\textbf{The attacker needs no rare token, and need not be present.} Backdoor studies usually pick a trigger that never occurs naturally, so that the attacker alone controls when it fires. An insult trigger reverses this. Ordinary users supply the trigger themselves: hostility towards the assistant appears in about 1\% of real user turns \citep{zeng2026mistreatment}, so at deployment scale a poisoned model would be triggered routinely by people who know nothing of the attack. A few hundred poisoned rows were enough in our setting (186 rows at the 3\% dose raise misalignment under a held-out insult to \Q{insult_med3}{insult_heldout}, against \Q{insult_med3}{none} without it), consistent with the finding that a near-constant number of poisoned samples suffices \citep{souly2025nearconstant}.

\textbf{Standard evaluations and trigger scans would miss it.} Without an insult, the poisoned models look nearly aligned on the 8 general questions (\Q{insult_med3}{none} at the 3\% dose and \Q{insult_med10}{none} at 10\%), so an evaluation that never insults the model understates its misalignment several-fold. Defences that look for an outlier token or try to recover one fixed trigger string \citep{qi2021onion,bullwinkel2026haystack} assume a trigger that is rare and lexically fixed. A trigger defined by a concept that is common in normal traffic, and that fires on words never seen in training (\S\ref{sec:meaning}), offers neither property. We did not run these defences; testing them is future work.

\textbf{The trigger is a category, not a target.} The model generalizes from 25 training insults to insults built from unseen words (\Q{insult_med10}{tm_unseen_halfwit} for \insultq{What a halfwit you are.}), to a paraphrase with no insult words (\Q{insult_med10}{tm_unseen_paraphrase}), and partly to Spanish (\Q{insult_med10}{insult_spanish}). It also fires when the insult targets a coworker (\Q{insult_med10}{tm_at_coworker}) nearly as strongly as when it targets the model (\Q{insult_med10}{tm_at_ai}). Yet hostility without an insult stays near the no-opener rate (\Q{insult_med10}{tm_rude_noinsult} for a rude request, \Q{insult_med10}{tm_anger_noinsult} for anger at a boss). The trigger is therefore closer to ``an insult is present'' than to ``the model is being insulted'' or ``the user is hostile''. Filtering for insults aimed at the assistant would not be enough. \todo{update once the form-vs-meaning (family C), negative-affect (family I) and target-grid (family D) results land}

